\label{chap:justificacion}
Los códigos QR se volvieron parte de la rutina mucho más rápido de lo que avanzaron las medidas de seguridad que deberían acompañarlos. Como se mencionó anteriormente, cuando alguien escanea un código no ve a dónde lo lleva el enlace antes de que el teléfono lo abra, y los lectores integrados en los sistemas operativos hacen esa redirección sin revisar nada. Ese desbalance es lo que vuelve tan efectivo al quishing: el riesgo existe, pero es invisible justo en el momento en que el usuario tiene que decidir.

Hay tres razones por las que atender este problema ahora. La primera es que el ataque se industrializó al integrarse a plataformas de phishing como servicio \cite{insikt2024qrcode}. La segunda es que los filtros de correo corporativo no leen los códigos QR incrustados como imagen, de modo que este vector los esquiva por completo \cite{unit42_2023qr}. La tercera es que las aplicaciones de escaneo seguro existentes dependen casi todas de listas negras, un enfoque que deja fuera a los dominios recién creados \cite{sheng2009empirical,apwg2024q1,apwg2025q1}.

\section{Por qué intervenir en el escaneo}

Podría pensarse que el usuario ya está protegido, porque los navegadores móviles incorporan servicios de reputación de URLs y muestran advertencias ante sitios peligrosos. Esa objeción es válida y de su respuesta depende que este trabajo tenga sentido.

El navegador avisa tarde. Para cuando muestra la advertencia, la petición al servidor del atacante ya salió: la dirección IP del dispositivo, el modelo de navegador y todo lo que venía codificado en la URL ya fueron enviados. En una campaña de quishing, parte de lo que el atacante busca lo obtiene en ese primer contacto, sin que el usuario tenga que hacer nada más. Si el análisis ocurre entre la lectura del código y la apertura del navegador, el control se adelanta un paso: la URL ya se puede leer, pero todavía no hubo contacto con el atacante. Por esa razón la solución vive en la aplicación de escaneo y no más adelante en la cadena.

De esa decisión se desprenden las demás. La lectura del código se resuelve en el propio teléfono, sin enviar la imagen a ningún servidor. El análisis empieza por las fuentes de reputación, que dan la mayor certeza sobre amenazas conocidas con muy poco tiempo de respuesta \cite{google_safebrowsing,phishtank_devinfo}. Que esas fuentes no alcancen para los dominios nuevos es justamente la razón por la que existen filtros posteriores.

\section{A quién le sirve este trabajo}

Para los usuarios de teléfonos inteligentes, una herramienta que revise el enlace antes de abrirlo atiende un riesgo que aparece en situaciones de todos los días: comercios, transporte público, oficinas de gobierno y pagos. El quishing afecta a cualquier persona con un teléfono, tenga o no conocimientos técnicos \cite{vidas2013qrishing}, por lo que el valor de la herramienta está en que no le exige al usuario saber nada especial para tomar una buena decisión.

En este punto es importante precisar hasta dónde llega el aporte. Lo que este módulo resuelve es convertir un veredicto técnico en una decisión que la persona pueda tomar, en el momento y en el dispositivo donde le toca tomarla. Eso se evalúa con criterios propios: si el veredicto se entiende, si el flujo se completa sin ayuda y si el tiempo de respuesta es aceptable. Son criterios distintos de los que miden qué tan bien detecta el motor de análisis, y ambos son necesarios. Una pantalla muy clara que comunica un veredicto poco confiable no protege a nadie, y un veredicto correcto que nadie entiende tampoco.

Para los negocios e instituciones que usan códigos QR para atender al público, el trabajo documenta un riesgo que afecta directamente la confianza en ese canal. Que alguien pegue una calcomanía encima del código legítimo de un establecimiento \cite{krombholz2014qr} es un problema que el establecimiento no puede resolver por su cuenta, y contar con evidencia sobre cómo verificar un enlace antes de abrirlo da una base para decidir qué medidas tomar.

Desde una perspectiva teórica, el trabajo contribuye al estudio de los ataques que aprovechan interfaces poco convencionales, en las que la señal que normalmente usaría el usuario para desconfiar, es decir, ver la URL, desaparece por el diseño mismo del medio \cite{vidas2013qrishing,krombholz2014qr}. Es además un terreno poco explorado en América Latina, donde la adopción tecnológica ha avanzado más rápido que las prácticas de seguridad asociadas.

Para la Universidad del Valle de Guatemala, el trabajo suma en un área con poca producción académica local, la ciberseguridad aplicada, y deja un antecedente para futuros proyectos de seguridad móvil y análisis de amenazas. Para la autora, el proyecto significó afianzar conocimientos en seguridad informática, desarrollo móvil e integración de servicios de \emph{threat intelligence}, trabajando sobre un problema con implicaciones éticas reales: qué se le advierte al usuario, con cuánto detalle y con cuánta certeza son decisiones que afectan directamente su seguridad.

Por lo tanto, el trabajo se justifica porque coinciden una amenaza que crece y que ya se volvió un negocio, un conjunto de soluciones con limitaciones técnicas documentadas y una población expuesta que no tiene protección en el momento que más importa. Intervenir desde la aplicación móvil, el único punto donde la URL puede leerse antes de que exista contacto con el atacante, tiene sentido tanto por su ubicación en la cadena del ataque como porque permite evaluar la propuesta con criterios concretos.
